% from ega3-1-fr.tex
\oldpage[0\textsubscript{III}]{28}
objets gradués de $K^\bullet$, autrement dit
$F^p(K^\bullet)=(K^i\cap F^p(K^\bullet))_{i\in\mathbb Z}$~; en outre, on
suppose que $d(F^p(K^\bullet))\subset F^p(K^\bullet)$ pour tout
$p\in\mathbb Z$. Rappelons alors rapidement comment on définit
\emph{fonctoriellement} une suite spectrale $E(K^\bullet)$ à partir de
$K^\bullet$ (M, XV, 4 et G, I, 4.3). Pour $r\geq2$, le morphisme canonique
$F^p(K^\bullet)/F^{p+r}(K^\bullet)\to
F^p(K^\bullet)/F^{p+1}(K^\bullet)$ définit un morphisme pour la cohomologie
\[
  \mathrm H^{p+q}(F^p(K^\bullet)/F^{p+r}(K^\bullet))
  \longrightarrow
  \mathrm H^{p+q}(F^p(K^\bullet)/F^{p+1}(K^\bullet)).
\]
On désigne par $Z_r^{pq}(K^\bullet)$ l'image de ce morphisme. De même, de la
suite exacte
\[
  0\to F^p(K^\bullet)/F^{p+1}(K^\bullet)
  \to F^{p-r+1}(K^\bullet)/F^{p+1}(K^\bullet)
  \to F^{p-r+1}(K^\bullet)/F^p(K^\bullet)\to0
\]
on déduit par la suite exacte de cohomologie un morphisme
\[
  \mathrm H^{p+q-1}(F^{p-r+1}(K^\bullet)/F^p(K^\bullet))
  \longrightarrow
  \mathrm H^{p+q}(F^p(K^\bullet)/F^{p+1}(K^\bullet))
\]
et on désigne par $B_r^{pq}(K^\bullet)$ l'image de ce morphisme~; on montre
que $B_r^{pq}(K^\bullet)\subset Z_r^{pq}(K^\bullet)$ et on prend
$E_r^{pq}(K^\bullet)=Z_r^{pq}(K^\bullet)/B_r^{pq}(K^\bullet)$~; nous ne
préciserons pas la définition des $d_r^{pq}$, ni des $\alpha_r^{pq}$.

On notera ici que tous les $Z_r^{pq}(K^\bullet)$ et
$B_r^{pq}(K^\bullet)$, pour $p$, $q$ fixés, sont des sous-objets du même
objet
$\mathrm H^{p+q}(F^p(K^\bullet)/F^{p+1}(K^\bullet))$, que l'on note
$Z_1^{pq}(K^\bullet)$~; on pose $B_1^{pq}(K^\bullet)=0$, de sorte que les
définitions précédentes pour $Z_r^{pq}(K^\bullet)$ et
$B_r^{pq}(K^\bullet)$ s'appliquent aussi pour $r=1$~; on pose encore
$E_1^{pq}(K^\bullet)=Z_1^{pq}(K^\bullet)$. On définit encore les $d_1^{pq}$
et $\alpha_1^{pq}$ de sorte que les conditions de
(\hyperref[0.11.1.1-fr]{11.1.1}) soient vérifiées pour $r=1$. On définit
d'autre part les sous-objets $Z_\infty^{pq}(K^\bullet)$, image du morphisme
\[
  \mathrm H^{p+q}(F^p(K^\bullet))\longrightarrow
  \mathrm H^{p+q}(F^p(K^\bullet)/F^{p+1}(K^\bullet))
  =E_1^{pq}(K^\bullet)
\]
et $B_\infty^{pq}(K^\bullet)$, image du morphisme
\[
  \mathrm H^{p+q-1}(K^\bullet/F^p(K^\bullet))\longrightarrow
  \mathrm H^{p+q}(F^p(K^\bullet)/F^{p+1}(K^\bullet))
  =E_1^{pq}(K^\bullet)
\]
déduit comme ci-dessus d'une suite exacte de cohomologie. On prend pour
$Z_\infty(E_2^{pq}(K^\bullet))$ et $B_\infty(E_2^{pq}(K^\bullet))$ les
images canoniques dans $E_2^{pq}(K^\bullet)$ de
$Z_\infty^{pq}(K^\bullet)$ et $B_\infty^{pq}(K^\bullet)$.

Enfin, on désigne par $F^p(\mathrm H^n(K^\bullet))$ l'image dans
$\mathrm H^n(K^\bullet)$ du morphisme
$\mathrm H^n(F^p(K^\bullet))\to\mathrm H^n(K^\bullet)$ provenant de
l'injection canonique $F^p(K^\bullet)\to K^\bullet$~; par la suite exacte de
cohomologie, c'est aussi le noyau du morphisme
$\mathrm H^n(K^\bullet)\to\mathrm H^n(K^\bullet/F^p(K^\bullet))$. On
définit ainsi une filtration sur
$E^n(K^\bullet)=\mathrm H^n(K^\bullet)$~; nous ne donnerons pas non plus ici
la définition des isomorphismes $\beta^{pq}$.
\end{env}

\begin{env}[11.2.3]
\label{0.11.2.3-fr}
Le caractère \emph{fonctoriel} de $E(K^\bullet)$ doit s'entendre de la façon
suivante~: étant donné deux complexes \emph{filtrés} $K^\bullet$,
$K^{\prime\bullet}$ de $\mathcal C$, et un morphisme de complexes
$u:K^\bullet\to K^{\prime\bullet}$ \emph{compatible avec les filtrations}, on
en déduit de façon évidente les morphismes $u_r^{pq}$ (pour $r\geq1$) et
$u^n$, et on montre que ces morphismes sont compatibles avec les $d_r^{pq}$,
$\alpha_r^{pq}$ et $\beta^{pq}$ au sens de
(\hyperref[0.11.1.2-fr]{11.1.2}), donc définissent bien un morphisme
$E(u):E(K^\bullet)\to E(K^{\prime\bullet})$ de suites spectrales. En outre,
on montre que si $u$ et $v$ sont des morphismes
$K^\bullet\to K^{\prime\bullet}$ du type précédent, \emph{homotopes d'ordre
$\leq k$}, alors $u_r^{pq}=v_r^{pq}$ pour $r>k$ et $u^n=v^n$ pour tout $n$
(M, XV, 3.1).
\end{env}

\begin{env}[11.2.4]
\label{0.11.2.4-fr}
